% LuaLaTeX document - compile with: luatex-pdf vce_architecture.tex
\documentclass[10pt,a4paper,twocolumn]{ltjsarticle}

% ファイル生成日時（JST）
\newcommand{\generatedDate}{2026-01-12}
\newcommand{\generatedTime}{19:00}

% 基本パッケージ
\usepackage{geometry}
\geometry{left=15mm,right=15mm,top=20mm,bottom=20mm}

% LuaLaTeX用フォント設定パッケージ
\usepackage{luatexja-fontspec}
\usepackage{amsmath,amssymb}
\usepackage{unicode-math}

% 欧文フォント設定 (Libertinus)
\setmainfont{Libertinus Serif}[
    BoldFont = {Libertinus Serif Bold},
    ItalicFont = {Libertinus Serif Italic},
    BoldItalicFont = {Libertinus Serif Bold Italic}
]
\setsansfont{Libertinus Sans}[
    BoldFont = {Libertinus Sans Bold},
    ItalicFont = {Libertinus Sans Italic}
]
\setmonofont{DejaVu Sans Mono}[Scale=0.8]

% 日本語フォント設定 (原ノ味フォント)
\setmainjfont{HaranoAjiMincho-Regular}[
    BoldFont = {HaranoAjiGothic-Medium},
    ItalicFont = {HaranoAjiMincho-Regular},
    BoldItalicFont = {HaranoAjiGothic-Bold}
]
\setsansjfont{HaranoAjiGothic-Regular}[
    BoldFont = {HaranoAjiGothic-Bold}
]
\setmonojfont{HaranoAjiGothic-Regular}

% 数式フォント設定 (Libertinus Math)
\setmathfont{Libertinus Math}

% ヘッダー・フッター設定
\usepackage{fancyhdr}
\usepackage{lastpage}
\pagestyle{fancy}
\fancyhf{}
\fancyhead[R]{\small \generatedDate\ \generatedTime\ JST (\thepage/\pageref{LastPage})}
\renewcommand{\headrulewidth}{0.4pt}

% 1ページ目のスタイル
\fancypagestyle{firstpage}{
    \fancyhf{}
    \renewcommand{\headrulewidth}{0pt}
}

% その他パッケージ
\usepackage{booktabs}
\usepackage{array}
\usepackage{tabularx}
\newcolumntype{Y}{>{\raggedright\arraybackslash}X}
\usepackage{ascmac}
\usepackage{hyperref}
\hypersetup{
    colorlinks=true,
    linkcolor=blue,
    urlcolor=blue
}
\usepackage{listings}
\lstset{
    basicstyle=\ttfamily\tiny,
    breaklines=true,
    frame=single,
    backgroundcolor=\color[gray]{0.95},
    xleftmargin=1mm,
    xrightmargin=1mm
}
\usepackage{xcolor}
\usepackage{enumitem}
\setlist{nosep,leftmargin=*}

% タイトル設定
\title{\textbf{VCE アーキテクチャ設計書}\\
{\large IPO・状態遷移・目的手段階層に基づく設計}}
\author{mashi--Claude Dialogue}
\date{}

\begin{document}
\maketitle
\thispagestyle{firstpage}

\section{設計手法}

本設計は以下の手法を組み合わせて導出された。

\vspace{0.5\baselineskip}
\noindent{\footnotesize
\begin{tabularx}{\linewidth}{@{}lY@{}}
\toprule
手法 & 適用箇所 \\
\midrule
IDEF0 (SADT) & IPOモデル、入力→中間→出力 \\
Goal-Oriented RE & 目的-手段の階層分解 \\
状態遷移モデル & ファイル状態の変化 \\
Responsibility-Driven & クラス境界と責務の決定 \\
\bottomrule
\end{tabularx}
}

\section{コンテンツフロー（IPO）}

\subsection{入力コンテンツ}

\vspace{0.3\baselineskip}
\noindent{\scriptsize
\begin{tabularx}{\linewidth}{@{}lYY@{}}
\toprule
コンテンツ & 形式 & 取得元 \\
\midrule
ソースメディア & MP4, MP3 & ローカル \\
YouTube動画 & URL→MP4 & yt-dlp \\
既存チャプター & .txt & 読込 \\
既存プロジェクト & .vce.json & 読込 \\
\bottomrule
\end{tabularx}
}

\subsection{中間生成物}

\subsubsection{永続化（人間の判断結果）}

\vspace{0.3\baselineskip}
\noindent{\scriptsize
\begin{tabularx}{\linewidth}{@{}lY@{}}
\toprule
コンテンツ & 説明 \\
\midrule
.vce.json & ソース・チャプター・設定を保持 \\
\bottomrule
\end{tabularx}
}

\subsubsection{導出構造（Derived）}

\vspace{0.3\baselineskip}
\noindent{\scriptsize
\begin{tabularx}{\linewidth}{@{}lYY@{}}
\toprule
構造 & 導出元 & 用途 \\
\midrule
VirtualTimeline & List[SourceFile] & 統一座標系 \\
ExcludedRegions & List[ChapterData] & ハッチング、カット \\
ExtractionPlan & Sources+Chapters & ffmpeg計画 \\
\bottomrule
\end{tabularx}
}

\subsubsection{分析結果（Analysis）}

\vspace{0.3\baselineskip}
\noindent{\scriptsize
\begin{tabularx}{\linewidth}{@{}lYY@{}}
\toprule
データ & 生成元 & 消費先 \\
\midrule
WaveformData & WaveformWorker & WaveformWidget \\
SpectrogramData & SpectroWorker & WaveformWidget \\
\bottomrule
\end{tabularx}
}

\subsection{最終出力}

\vspace{0.3\baselineskip}
\noindent{\scriptsize
\begin{tabularx}{\linewidth}{@{}lYY@{}}
\toprule
コンテンツ & 形式 & 生成元 \\
\midrule
エンコード済み & MP4(章埋込) & vce-encode \\
分割ファイル群 & MP4/MP3 & vce-split \\
チャプター & .txt & vce-chapters \\
\bottomrule
\end{tabularx}
}

\section{状態遷移モデル}

\subsection{状態定義}

\vspace{0.3\baselineskip}
\noindent{\scriptsize
\begin{tabularx}{\linewidth}{@{}clYY@{}}
\toprule
状態 & 名称 & データ構造 & 説明 \\
\midrule
S0 & 外部 & Path & ファイル位置 \\
S1 & 読込済 & SourceFile & パス+メタ \\
S1' & 座標系 & VirtualTimeline & 統一座標（導出） \\
S2 & 分析済 & Waveform/Spectro & 可視化用 \\
S3 & 構造化 & ChapterData & 位置+名前+除外 \\
S3' & 除外計算 & ExcludedRegions & 除外区間（導出） \\
S4 & 永続化 & VCEProject & 全データ保持 \\
S4' & 計画 & ExtractionPlan & セグメント（導出） \\
S5 & 出力 & ExportResult & 成否+パス \\
\bottomrule
\end{tabularx}
}

\subsection{状態遷移図}

\begin{lstlisting}
[S0] --取得--> [S1] --導出--> [S1']
                |               |
                | 波形抽出      | 座標系提供
                v               v
              [S2] <------------+
                |
                | 人間判断(UIロック)
                v
              [S3] --導出--> [S3']
                |               |
                | 保存          |
                v               v
              [S4] --導出--> [S4']
                |               |
                | エンコード    | セグメント情報
                v               v
              [S5] <------------+
\end{lstlisting}

\subsection{因果関係マトリックス}

\vspace{0.3\baselineskip}
\noindent{\tiny
\begin{tabularx}{\linewidth}{@{}lYYYYYYYY@{}}
\toprule
 & S0→1 & 1→1' & 1→2 & 2→3 & 3→3' & 3→4 & 4→4' & 4→5 \\
\midrule
取得 & ● & & & & & & & \\
座標系 & 前 & ● & & & & & & \\
波形 & 前 & 前 & ● & & & & & \\
判断 & & 参 & 参 & ● & & & & \\
除外 & & & & 前 & ● & & & \\
永続 & & & & 前 & & ● & & \\
計画 & & & & & 前 & 前 & ● & \\
出力 & & & & & & & 前 & ● \\
\bottomrule
\end{tabularx}
}

{\tiny 凡例: ●直接 / 前=前提 / 参=参照}

\section{目的-手段階層}

\subsection{レベル0: 最上位目的}

\begin{itembox}[l]{P0}
素材を構造化された再利用可能な形式に変換する
\end{itembox}

\subsection{レベル1: 目的分解}

\vspace{0.3\baselineskip}
\noindent{\scriptsize
\begin{tabularx}{\linewidth}{@{}clY@{}}
\toprule
ID & 目的 & 説明 \\
\midrule
P1.1 & 素材取得 & 外部メディアを取込 \\
P1.2 & 判断支援 & 人間が判断可能に \\
P1.3 & 構造定義 & 時間軸構造を定義 \\
P1.4 & 永続化 & 判断結果を保存 \\
P1.5 & 成果物 & 最終出力を生成 \\
\bottomrule
\end{tabularx}
}

\subsection{レベル2: 手段}

\vspace{0.3\baselineskip}
\noindent{\scriptsize
\begin{tabularx}{\linewidth}{@{}clY@{}}
\toprule
目的 & ID & 手段 \\
\midrule
P1.1 & M1.1.1 & ファイル読込 \\
 & M1.1.2 & YouTube DL \\
 & M1.1.3 & メタデータ検出 \\
 & M1.1.4 & 仮想TL構築 \\
\midrule
P1.2 & M1.2.1 & 動画プレビュー \\
 & M1.2.2 & 音声再生 \\
 & M1.2.3 & 波形表示 \\
 & M1.2.4 & スペクトログラム \\
 & M1.2.5 & ナビゲーション \\
\midrule
P1.3 & M1.3.1 & チャプター追加 \\
 & M1.3.2 & 除外区間設定 \\
 & M1.3.3 & ソース順序変更 \\
 & M1.3.4 & ハッチング表示 \\
\midrule
P1.4 & M1.4.1 & .vce.json書出 \\
 & M1.4.2 & チャプター書出 \\
\midrule
P1.5 & M1.5.1 & カット計画 \\
 & M1.5.2 & エンコード \\
 & M1.5.3 & チャプター埋込 \\
 & M1.5.4 & タイトル焼込 \\
 & M1.5.5 & 分割出力 \\
\bottomrule
\end{tabularx}
}

\section{データ構造}

\subsection{基本データ}

\begin{lstlisting}[language=Python]
@dataclass
class SourceFile:
    path: Path
    duration_ms: int
    media_type: Literal["video","audio"]
    metadata: MediaInfo

@dataclass
class MediaInfo:
    codec: str
    resolution: Optional[Tuple[int,int]]
    frame_rate: Optional[float]
    audio_channels: int
    audio_sample_rate: int
\end{lstlisting}

\subsection{導出構造}

\begin{lstlisting}[language=Python]
@dataclass
class VirtualTimeline:
    """複数ソースを統一座標系に射影"""
    total_duration_ms: int
    source_boundaries: List[int]
    source_offsets: List[int]

    def virtual_to_source(self, ms)
        -> Tuple[int, int]: ...
    def source_to_virtual(self, idx, ms)
        -> int: ...

@dataclass
class ExcludedRegion:
    start_ms: int
    end_ms: int

@dataclass
class SegmentInfo:
    source_index: int
    start_ms: int
    end_ms: int

@dataclass
class ExtractionPlan:
    segments: List[SegmentInfo]
    total_output_duration_ms: int
\end{lstlisting}

\subsection{分析結果}

\begin{lstlisting}[language=Python]
@dataclass
class WaveformData:
    source_index: int
    samples: np.ndarray
    sample_rate: int
    duration_ms: int

@dataclass
class SpectrogramData:
    source_index: int
    image: np.ndarray  # 2D (freq x time)
    freq_range: Tuple[int, int]
    duration_ms: int
\end{lstlisting}

\subsection{判断結果・永続化}

\begin{lstlisting}[language=Python]
@dataclass
class ChapterData:
    position_ms: int  # 仮想TL座標
    title: str
    is_excluded: bool = False

@dataclass
class ExportSettings:
    encoder: str
    quality: int  # 0-5
    overlay_title: bool
    embed_chapters: bool

@dataclass
class VCEProject:
    version: str
    sources: List[SourceFile]
    chapters: List[ChapterData]
    export_settings: ExportSettings
    created_at: datetime
    modified_at: datetime
\end{lstlisting}

\subsection{出力・状態}

\begin{lstlisting}[language=Python]
@dataclass
class ExportResult:
    success: bool
    output_path: Path
    duration_ms: int
    error_message: Optional[str]

class ExportState(Enum):
    IDLE = "idle"
    EXTRACTING = "extracting"
    ENCODING = "encoding"
    EMBEDDING = "embedding"
    COMPLETED = "completed"
    ERROR = "error"
    CANCELLED = "cancelled"

@dataclass
class ProjectState:
    is_dirty: bool
    last_saved_at: Optional[datetime]
\end{lstlisting}

\section{クラス設計}

\subsection{目的→クラス写像}

\vspace{0.3\baselineskip}
\noindent{\scriptsize
\begin{tabularx}{\linewidth}{@{}clY@{}}
\toprule
目的 & クラス & 責務 \\
\midrule
P1.1 & SourceFileMgr & S0→S1 \\
P1.1 & PlaybackMgr & S1→S1' VirtualTL \\
P1.2 & WaveformWorker & S1→S2 波形 \\
P1.2 & SpectroWorker & S1→S2 スペクトロ \\
P1.2 & WaveformWidget & 表示 \\
P1.3 & ChapterMgr & S2→S3, S3→S3' \\
P1.4 & ProjectPers & S3↔S4 \\
P1.5 & ExportOrch & S4→S4'→S5 \\
\bottomrule
\end{tabularx}
}

\subsection{SourceFileManager}

\begin{lstlisting}
責務: S0 -> S1 遷移管理

入力: Path, URL
出力: List[SourceFile], シグナル

Interface:
  + add_source(path) -> SourceFile
  + add_from_url(url) -> SourceFile
  + remove_source(index)
  + reorder_sources(new_order)
  + get_sources() -> List[SourceFile]

依存: なし（最上流）
\end{lstlisting}

\subsection{PlaybackManager}

\begin{lstlisting}
責務:
  1. S1->S1' VirtualTimeline構築
  2. 再生制御
  3. 座標変換

保持: VirtualTimeline

Interface:
  + set_sources(sources)
  + get_virtual_timeline() -> VirtualTimeline
  + play() / pause() / seek(ms)
  + get_position() -> int
  + virtual_to_source(ms) -> (idx, ms)
  + source_to_virtual(idx, ms) -> ms
  + get_total_duration() -> int
  + get_source_boundaries() -> List[int]

依存: SourceFileManager
\end{lstlisting}

\subsection{WaveformWorker}

\begin{lstlisting}
責務: S1 -> S2 分析データ生成

入力: SourceFile
出力: WaveformData/SpectrogramData

Interface:
  + start(source: SourceFile)
  + cancel()
  シグナル: progress(%), finished(data)

依存: SourceFile
\end{lstlisting}

\subsection{ChapterManager}

\begin{lstlisting}
責務:
  1. S2->S3 チャプター管理
  2. S3->S3' ExcludedRegions導出

保持: List[ChapterData]

Interface:
  + add_chapter(pos, title) -> ChapterData
  + remove_chapter(index)
  + update_chapter(index, **kwargs)
  + set_excluded(index, bool)
  + get_chapters() -> List[ChapterData]
  + get_excluded_regions()
      -> List[ExcludedRegion]

依存: PlaybackManager (VirtualTimeline)
\end{lstlisting}

\subsection{ProjectPersistence}

\begin{lstlisting}
責務: S3 <-> S4 永続化/読込

Interface:
  + save_project(path, project)
  + load_project(path) -> VCEProject
  + export_chapters(path, chapters)
  + import_chapters(path) -> List

依存: なし（データ変換のみ）
\end{lstlisting}

\subsection{ExportOrchestrator}

\begin{lstlisting}
責務:
  1. S4->S4' ExtractionPlan導出
  2. S4'->S5 エクスポート実行

保持: ExportState, ExtractionPlan

Interface:
  + calculate_extraction_plan(project)
      -> ExtractionPlan
  + start_export(project, settings)
  + start_split(project, settings)
  + cancel()
  + get_state() -> ExportState
  シグナル: progress(%), state_changed,
            export_completed

依存: VCEProject, ExcludedRegions
\end{lstlisting}

\section{依存関係図}

\begin{lstlisting}
SourceFileManager (S0->S1)
        |
        | List[SourceFile]
        v
PlaybackManager (S1->S1')
   [VirtualTimeline]
        |
   +----+----+
   |    |    |
   v    v    v
Waveform Spectro (座標系参照)
Worker   Worker      |
   |       |         |
   v       v         v
WaveformWidget  ChapterManager
(判断支援)       (S2->S3, S3->S3')
  - 波形描画     [ExcludedRegions]
  - スペクトロ        |
  - ハッチング        |
  - 境界線     <------+
                      |
                      v
              ProjectPersistence
                  (S3<->S4)
                      |
                      v
              ExportOrchestrator
               (S4->S4'->S5)
              [ExtractionPlan]
              [ExportState]
                      |
     +----------------+----------------+
     v                v                v
vce-encode       vce-split       vce-chapters
  (CLI)            (CLI)            (CLI)
\end{lstlisting}

\section{トレーサビリティ}

\vspace{0.3\baselineskip}
\noindent{\tiny
\begin{tabularx}{\linewidth}{@{}llYY@{}}
\toprule
Why & How & Who & What \\
\midrule
P1.1 & M1.1.1 & SrcFileMgr & SourceFile \\
 & M1.1.3 & & MediaInfo \\
 & M1.1.4 & PlaybackMgr & VirtualTimeline \\
\midrule
P1.2 & M1.2.3 & WaveformWkr & WaveformData \\
 & M1.2.4 & SpectroWkr & SpectrogramData \\
\midrule
P1.3 & M1.3.1 & ChapterMgr & ChapterData \\
 & M1.3.2 & & ExcludedRegion \\
\midrule
P1.4 & M1.4.1 & ProjectPers & VCEProject \\
\midrule
P1.5 & M1.5.1 & ExportOrch & ExtractionPlan \\
 & M1.5.2 & & ExportState \\
 & M1.5.5 & & ExportResult \\
\bottomrule
\end{tabularx}
}

\section{作用の分類}

\subsection{直接作用（状態遷移）}

\vspace{0.3\baselineskip}
\noindent{\tiny
\begin{tabularx}{\linewidth}{@{}lYYY@{}}
\toprule
作用 & 遷移 & クラス & データ \\
\midrule
ファイル取得 & S0→S1 & SrcFileMgr & SourceFile \\
座標系構築 & S1→S1' & PlaybackMgr & VirtualTL \\
波形抽出 & S1→S2 & WaveformWkr & WaveformData \\
スペクトロ & S1→S2 & SpectroWkr & SpectroData \\
チャプター & S2→S3 & ChapterMgr & ChapterData \\
除外計算 & S3→S3' & ChapterMgr & ExcludedRgn \\
永続化 & S3→S4 & ProjectPers & VCEProject \\
計画立案 & S4→S4' & ExportOrch & ExtractPlan \\
エンコード & S4'→S5 & ExportOrch & ExportResult \\
\bottomrule
\end{tabularx}
}

\subsection{間接作用（支援）}

\vspace{0.3\baselineskip}
\noindent{\scriptsize
\begin{tabularx}{\linewidth}{@{}lYY@{}}
\toprule
作用 & 支援対象 & クラス \\
\midrule
波形表示 & 人間判断→チャプ & WaveformWgt \\
スペクトロ & 人間判断→チャプ & WaveformWgt \\
プレビュー & 人間判断→チャプ & PlaybackMgr \\
ハッチング & 判断→除外確認 & WaveformWgt \\
境界線 & 判断→ソース確認 & WaveformWgt \\
\bottomrule
\end{tabularx}
}

\subsection{制約作用（前提条件）}

\vspace{0.3\baselineskip}
\noindent{\scriptsize
\begin{tabularx}{\linewidth}{@{}lY@{}}
\toprule
制約 & 内容 \\
\midrule
S1必須 & VirtualTimeline構築の前提 \\
S1'必須 & 波形抽出は座標系確定後 \\
S2必須 & 人間判断には分析データが必要 \\
S3必須 & 永続化には構造化が前提 \\
S3'必須 & 計画立案には除外計算が前提 \\
S4'必須 & エンコードには計画が前提 \\
\bottomrule
\end{tabularx}
}

\section*{参照ドキュメント}

\begin{itemize}
\item DESIGN\_PRINCIPLES.md -- 設計原則
\item vce\_feature\_matrix.md -- 機能マトリックス
\end{itemize}

\end{document}
